%% schriftsatz.latex — Anpassbares Pandoc-Template für familienrechtliche Schriftsätze
%%
%% Pandoc-Variablen (via -V oder YAML-Frontmatter):
%%   mainfont      Schriftart             (Standard: Arial)
%%   fontsize      Schriftgröße           (Standard: 11pt)
%%   papersize     Papierformat           (Standard: a4)
%%   geometry      Seitenränder           (Standard: margin=2.5cm)
%%   linestretch   Zeilenabstand          (Standard: 1.2)
%%   lang          Sprache                (Standard: de-DE)
%%   briefkopf     Vorbereiteter LaTeX-Briefkopf (von generate-pdf.py erzeugt)
%%
%% ─────────────────────────────────────────────────────────────────────────────

\documentclass[
  11pt,
  a4paper,
]{article}

%% ── Sprache ──────────────────────────────────────────────────────────────────
\usepackage[ngerman]{babel}

%% ── Fonts (XeLaTeX) ──────────────────────────────────────────────────────────
\usepackage{fontspec}
\setmainfont{Arial}[
  BoldFont       = Arial Bold,
  ItalicFont     = Arial Italic,
  BoldItalicFont = Arial Bold Italic,
  Ligatures      = TeX,
]
\usepackage{microtype}

%% ── Seitenlayout ─────────────────────────────────────────────────────────────
\usepackage[left=3cm,right=2cm,top=2.5cm,bottom=2.5cm]{geometry}

%% ── Zeilenabstand & Absätze ──────────────────────────────────────────────────
\usepackage{setspace}
\setstretch{1.3}
\usepackage{parskip}
\setlength{\parskip}{6pt}
\setlength{\parindent}{0pt}

%% ── Überschriften ────────────────────────────────────────────────────────────
\usepackage{titlesec}
\titleformat{\section}{\normalfont\large\bfseries}{}{0em}{}
\titleformat{\subsection}{\normalfont\normalsize\bfseries}{}{0em}{}
\titleformat{\subsubsection}{\normalfont\normalsize\bfseries}{}{0em}{}
\titlespacing{\section}{0pt}{14pt}{6pt}
\titlespacing{\subsection}{0pt}{10pt}{4pt}
%% Keine automatische Abschnittsnummerierung
\setcounter{secnumdepth}{0}

%% ── Seitenzahlen ─────────────────────────────────────────────────────────────
\usepackage{fancyhdr}
\usepackage{lastpage}
\pagestyle{fancy}
\fancyhf{}
\fancyfoot[C]{\footnotesize Seite~\thepage~von~\pageref{LastPage}}
\renewcommand{\headrulewidth}{0pt}
\renewcommand{\footrulewidth}{0pt}

%% ── Tabellen ─────────────────────────────────────────────────────────────────
\usepackage{booktabs}
\usepackage{longtable}
\usepackage{array}
\usepackage{calc}
\usepackage{caption}
\captionsetup{font=small, labelfont=bf}
%% Pandoc 3.x setzt \LTcaptype auf "none" — LaTeX braucht diesen Counter:
\newcounter{none}

%% ── Listen ───────────────────────────────────────────────────────────────────
\usepackage{enumitem}
\setlist{noitemsep, topsep=4pt}

%% ── Sonstiges ────────────────────────────────────────────────────────────────
\usepackage{xcolor}
\usepackage[normalem]{ulem}
\usepackage[hidelinks]{hyperref}
\usepackage{needspace}

%% ── Pandoc-Kompatibilität ────────────────────────────────────────────────────
\providecommand{\tightlist}{%
  \setlength{\itemsep}{0pt}\setlength{\parskip}{0pt}}

%% ─────────────────────────────────────────────────────────────────────────────
\begin{document}

%% Briefkopf — von generate-pdf.py aus Markdown-Zeilen vor dem ersten --- erzeugt.
%% Zeilen mit "den \d" werden rechtsbündig gesetzt, Betreff fett+unterstrichen,
%% Aktenzeichen fett, alle anderen Zeilen linksbündig mit explizitem Zeilenumbruch.

\section{Anlagenverzeichnis --- Az. 7 F
88/25}\label{anlagenverzeichnis-az.-7-f-8825}

{\def\LTcaptype{none} % do not increment counter
\begin{longtable}[]{@{}
  >{\raggedright\arraybackslash}p{(\linewidth - 10\tabcolsep) * \real{0.1250}}
  >{\raggedright\arraybackslash}p{(\linewidth - 10\tabcolsep) * \real{0.1094}}
  >{\raggedright\arraybackslash}p{(\linewidth - 10\tabcolsep) * \real{0.1094}}
  >{\raggedright\arraybackslash}p{(\linewidth - 10\tabcolsep) * \real{0.0781}}
  >{\raggedright\arraybackslash}p{(\linewidth - 10\tabcolsep) * \real{0.4531}}
  >{\raggedright\arraybackslash}p{(\linewidth - 10\tabcolsep) * \real{0.1250}}@{}}
\toprule\noalign{}
\begin{minipage}[b]{\linewidth}\raggedright
Anlage
\end{minipage} & \begin{minipage}[b]{\linewidth}\raggedright
Titel
\end{minipage} & \begin{minipage}[b]{\linewidth}\raggedright
Datei
\end{minipage} & \begin{minipage}[b]{\linewidth}\raggedright
Typ
\end{minipage} & \begin{minipage}[b]{\linewidth}\raggedright
Im Schriftsatz referenziert
\end{minipage} & \begin{minipage}[b]{\linewidth}\raggedright
Status
\end{minipage} \\
\midrule\noalign{}
\endhead
\bottomrule\noalign{}
\endlastfoot
B1 & Elternvereinbarung vom 01.09.2025 & elternvereinbarung-20250901.pdf
& Original & Ja --- Abschnitt: II.3, Anträge & beigefügt \\
B2 & KiTa-Bericht Regenbogen, Halbjahresbericht Herbst 2025 &
kita-bericht-herbst-2025.pdf & Kopie & Ja --- Abschnitt:
Glaubhaftmachung & beigefügt \\
B3 & WhatsApp-Verlauf AS/AG Oktober--November 2025 &
whatsapp-verlauf-okt-nov-2025.pdf & Kopie & Ja --- Abschnitt: VI,
Glaubhaftmachung & beigefügt \\
B4 & Dienstreise-Belege (Hotelrechnung Stuttgart + Zugtickets
Düsseldorf) & dienstreise-belege-2025-2026.pdf & Kopie & Ja ---
Abschnitt: II.2, VII, Glaubhaftmachung & beigefügt \\
B5 & E-Mail KiTa-Entwicklungsgespräch vom 10.02.2026 &
kita-email-20260210.pdf & Original & Ja --- Abschnitt: IV,
Glaubhaftmachung & beigefügt \\
B6 & Gemeinsamer Jahreskalender 2025/2026 (Google-Calendar-Export) &
jahreskalender-export-2026.pdf & Export & Ja --- Abschnitt: Anträge Nr.
3, Glaubhaftmachung & beigefügt \\
\end{longtable}
}

\subsection{Hinweise}\label{hinweise}

\begin{itemize}
\tightlist
\item
  \textbf{Original}: Wird physisch eingereicht. Der Skill generiert ein
  Deckblatt (deckblatt-B1.docx).
\item
  \textbf{Kopie}: Ausreichend als Ausdruck. Kein Deckblatt nötig.
\item
  \textbf{Export}: Aus Markdown generiert (Kalender, Protokoll). Kein
  separates Deckblatt.
\item
  Anlagen die bereits beim Gericht liegen, müssen nicht erneut
  eingereicht werden.
\end{itemize}

\end{document}
